\chapter{O Arquivo de Dados \texttt{ecommerce\_join.csv}}
\label{ap:csv}

Este apêndice reproduz, em forma tabular, o arquivo de dados do cenário exemplo,
\texttt{data/ecommerce/ecommerce\_join.csv}, utilizado nas seções~\ref{sec:algoritmo}
e~\ref{sec:exemplo-uso}. O arquivo contém 32 linhas de dados e 39 colunas, e resulta da
\textbf{junção} (\textit{join}) de todas as coleções de dados da fonte --- usuários,
endereços, produtos, categorias, estoque, pedidos, seleções de produto e carrinhos de
compra --- em uma única tabela larga. Por isso ele é \textbf{esparso}: cada linha
representa um evento da loja virtual, identificado pela coluna \texttt{action}
(\texttt{stock}, \texttt{purchase}, \texttt{select\_product} ou \texttt{add\_to\_cart}),
e preenche apenas as colunas pertinentes àquele tipo de evento --- as demais permanecem
vazias. A Tabela~\ref{tab:csv-colunas-acao} resume quais grupos de colunas cada ação
preenche. Este arquivo é a entrada do \textbf{modo combinado} de importação; no
\textbf{modo multifonte}, os mesmos dados são fornecidos como quatro arquivos CSV separados,
um por origem (\texttt{ecommerce\_stock.csv}, \texttt{ecommerce\_purchase.csv},
\texttt{ecommerce\_select\_product.csv} e \texttt{ecommerce\_add\_to\_cart.csv}).

\begin{table}[H]
  \caption{Grupos de colunas do \texttt{ecommerce\_join.csv} preenchidos por ação.}
  \label{tab:csv-colunas-acao}
  \centering
  {
  \setlength{\extrarowheight}{4pt}%
  \begin{tabular}{>{\raggedright\arraybackslash}p{2.6cm} >{\raggedright\arraybackslash}p{5.6cm} c c c c}
    \toprule
    \textbf{Grupo} & \textbf{Colunas} & \rotatebox{60}{\texttt{stock}} & \rotatebox{60}{\texttt{purchase}} & \rotatebox{60}{\texttt{select\_product}} & \rotatebox{60}{\texttt{add\_to\_cart}} \\
    \midrule
    Evento & \texttt{action}, \texttt{timestamp} & $\bullet$ & $\bullet$ & $\bullet$ & $\bullet$ \\
    Usuário & \texttt{user\_id}, \texttt{user\_name}, \texttt{user\_email} & $\bullet$ & $\bullet$ & $\bullet$ & $\bullet$ \\
    Endereço do usuário & \texttt{street}, \texttt{neighborhood}, \texttt{state}, \texttt{country}, \texttt{zip\_code} & $\bullet$ & $\bullet$ & --- & --- \\
    Pedido & \texttt{order\_number}, \texttt{order\_date}, \texttt{order\_status} & --- & $\bullet$ & --- & --- \\
    Contraparte da venda & \texttt{traded\_with}, \texttt{trader\_street}, \texttt{trader\_neighborhood}, \texttt{trader\_state}, \texttt{trader\_country}, \texttt{trader\_zip\_code} & --- & $\bullet$ & --- & --- \\
    Avaliação e pagamento & \texttt{comment}, \texttt{rating}, \texttt{payment\_method}, \texttt{quantity} & --- & $\bullet$ & --- & --- \\
    Produto & \texttt{product\_id}, \texttt{product\_name}, \texttt{product\_variant}, \texttt{product\_brand}, \texttt{product\_description}, \texttt{product\_image} & $\bullet$ & $\bullet$ & --- & --- \\
    Categoria & \texttt{category\_id}, \texttt{category\_name} & $\bullet$ & $\bullet$ & --- & --- \\
    Estoque & \texttt{quantity\_available}, \texttt{price}, \texttt{last\_restock\_date} & $\bullet$ & $\bullet$ & --- & --- \\
    Seleção de produto & \texttt{selected\_product\_id}, \texttt{suggested\_product\_count} & --- & --- & $\bullet$ & --- \\
    Carrinho & \texttt{shopping\_cart\_id}, \texttt{cart\_product\_id}, \texttt{cart\_quantity} & --- & --- & --- & $\bullet$ \\
    \bottomrule
  \end{tabular}}
  \par\vspace{4pt}
  {\footnotesize Fonte: Elaborado pelo autor (2026).}
\end{table}

Vários valores do arquivo seguem padrões regulares, o que permite apresentá-los de
forma resumida sem perda de informação. O produto \texttt{Product}$N$ tem
\texttt{product\_id} $N$, pertence à categoria \texttt{Categoria}$N$
(\texttt{category\_id} $N$) e tem descrição e imagem padronizadas
(\texttt{Description}$M$ e \texttt{https://example.com/image}$M$\texttt{.jpg}). Cada
usuário \texttt{user}$N$ tem e-mail \texttt{user}$N$\texttt{@example.com}, rua
\texttt{Street }$N$\texttt{, 42} e um identificador alfanumérico
(\texttt{user\_id}) cujo início é dado na Tabela~\ref{tab:csv-usuarios}, junto com os
demais dados de endereço. Nas tabelas deste apêndice, os usuários são referidos pelo
nome (\texttt{user}$N$), tanto na coluna \texttt{user\_id} quanto em
\texttt{traded\_with}.

\begin{table}[H]
  \caption{Usuários do \texttt{ecommerce\_join.csv} e seus dados de identificação e endereço.}
  \label{tab:csv-usuarios}
  \centering
  \begin{tabular}{l l l c l r}
    \toprule
    \textbf{Usuário} & \textbf{\texttt{user\_id} (início)} & \textbf{Bairro} & \textbf{Estado} & \textbf{País} & \textbf{CEP} \\
    \midrule
    user1 & \texttt{7fd7eee7-\ldots} & ABC & CA & United States & 12345 \\
    user2 & \texttt{f246fb91-\ldots} & DEF & SC & Brasil & 84102149 \\
    user3 & \texttt{09b3458a-\ldots} & GHI & RS & Brasil & 18057313 \\
    user4 & \texttt{1fd3d09a-\ldots} & JKL & TO & Brasil & 19645442 \\
    user5 & \texttt{0331463b-\ldots} & MNO & RJ & Brasil & 30097508 \\
    user6 & \texttt{f4857da9-\ldots} & PQR & SP & Brasil & 92083256 \\
    user7 & \texttt{a87df805-\ldots} & STU & SP & Brasil & 67065799 \\
    user8 & \texttt{c11a1ef2-\ldots} & VWX & SE & Brasil & 18439703 \\
    user9 & \texttt{d0416923-\ldots} & YZ  & AM & Brasil & 18439666 \\
    \bottomrule
  \end{tabular}
  \par\vspace{4pt}
  {\footnotesize Fonte: Elaborado pelo autor (2026).}
\end{table}

Nas seções seguintes, o arquivo é apresentado em quatro blocos, um por ação. As tabelas são
\textbf{verticais}: cada linha da tabela corresponde a uma coluna do CSV, e cada coluna da
tabela, a uma linha de dados daquela ação, na ordem em que aparece no arquivo. Todas as
colunas que a ação preenche são mostradas, com os valores exatamente como no arquivo,
exceto \texttt{user\_id} e \texttt{traded\_with}, que aparecem pelo nome do usuário
(Tabela~\ref{tab:csv-usuarios}). Como cada ação tem oito linhas, cada bloco ocupa duas
tabelas, com as linhas 1 a 4 e 5 a 8.

\section{Linhas com \texttt{action == stock}}

As oito linhas de reposição de estoque preenchem os grupos de evento, usuário, endereço,
produto, categoria e estoque (21 das 39 colunas). Todas ocorrem em 02/11/2023 e, em cada
uma, \texttt{last\_restock\_date} repete o instante do próprio evento
(\texttt{timestamp}). As Tabelas~\ref{tab:csv-stock-completa}
e~\ref{tab:csv-stock-completa-b} apresentam os valores.

\begin{table}[H]
  \caption{Linhas \texttt{action == stock} do \texttt{ecommerce\_join.csv}, todas as colunas preenchidas (linhas 1 a 4).}
  \label{tab:csv-stock-completa}
  \centering
  \setlength{\tabcolsep}{3pt}%
  \renewcommand{\arraystretch}{1.15}%
  \begin{tabular}{l *{4}{>{\raggedright\arraybackslash}p{2.5cm}}}
    \toprule
    \textbf{Coluna} & \textbf{Linha 1} & \textbf{Linha 2} & \textbf{Linha 3} & \textbf{Linha 4} \\
    \midrule
    \texttt{action} & stock & stock & stock & stock \\
    \texttt{timestamp} & 2023-11-02 03:30:00Z & 2023-11-02 06:45:00Z & 2023-11-02 09:15:00Z & 2023-11-02 12:00:00Z \\
    \texttt{user\_id} & user1 & user2 & user2 & user8 \\
    \texttt{user\_name} & user1 & user2 & user2 & user8 \\
    \texttt{user\_email} & user1@\allowbreak example.com & user2@\allowbreak example.com & user2@\allowbreak example.com & user8@\allowbreak example.com \\
    \texttt{street} & Street 1, 42 & Street 2, 42 & Street 2, 42 & Street 8, 42 \\
    \texttt{neighborhood} & ABC & DEF & DEF & VWX \\
    \texttt{state} & CA & SC & SC & SE \\
    \texttt{country} & United States & Brasil & Brasil & Brasil \\
    \texttt{zip\_code} & 12345 & 84102149 & 84102149 & 18439703 \\
    \texttt{product\_id} & 2 & 3 & 9 & 8 \\
    \texttt{product\_name} & Product2 & Product3 & Product9 & Product8 \\
    \texttt{product\_variant} & blue & green & black & white \\
    \texttt{product\_brand} & Microsoft & Google & Pepsi & Google \\
    \texttt{product\_description} & Description1 & Description2 & Description8 & Description7 \\
    \texttt{product\_image} & https://\allowbreak example.com/\allowbreak image1.jpg & https://\allowbreak example.com/\allowbreak image2.jpg & https://\allowbreak example.com/\allowbreak image8.jpg & https://\allowbreak example.com/\allowbreak image7.jpg \\
    \texttt{category\_id} & 2 & 3 & 9 & 8 \\
    \texttt{category\_name} & Categoria2 & Categoria3 & Categoria9 & Categoria8 \\
    \texttt{quantity\_available} & 2 & 1 & 67 & 60 \\
    \texttt{price} & 50.5 & 52.07 & 32.78 & 46.79 \\
    \texttt{last\_restock\_date} & 2023-11-02 03:30:00Z & 2023-11-02 06:45:00Z & 2023-11-02 09:15:00Z & 2023-11-02 12:00:00Z \\
    \bottomrule
  \end{tabular}
  \par\vspace{4pt}
  {\footnotesize Fonte: Elaborado pelo autor (2026), a partir de \texttt{data/ecommerce/ecommerce\_join.csv}.}
\end{table}

\begin{table}[H]
  \caption{Linhas \texttt{action == stock} do \texttt{ecommerce\_join.csv}, todas as colunas preenchidas (linhas 5 a 8).}
  \label{tab:csv-stock-completa-b}
  \centering
  \setlength{\tabcolsep}{3pt}%
  \renewcommand{\arraystretch}{1.15}%
  \begin{tabular}{l *{4}{>{\raggedright\arraybackslash}p{2.5cm}}}
    \toprule
    \textbf{Coluna} & \textbf{Linha 5} & \textbf{Linha 6} & \textbf{Linha 7} & \textbf{Linha 8} \\
    \midrule
    \texttt{action} & stock & stock & stock & stock \\
    \texttt{timestamp} & 2023-11-02 14:30:00Z & 2023-11-02 17:15:00Z & 2023-11-02 20:00:00Z & 2023-11-02 22:30:00Z \\
    \texttt{user\_id} & user9 & user9 & user9 & user9 \\
    \texttt{user\_name} & user9 & user9 & user9 & user9 \\
    \texttt{user\_email} & user9@\allowbreak example.com & user9@\allowbreak example.com & user9@\allowbreak example.com & user9@\allowbreak example.com \\
    \texttt{street} & Street 9, 42 & Street 9, 42 & Street 9, 42 & Street 9, 42 \\
    \texttt{neighborhood} & YZ & YZ & YZ & YZ \\
    \texttt{state} & AM & AM & AM & AM \\
    \texttt{country} & Brasil & Brasil & Brasil & Brasil \\
    \texttt{zip\_code} & 18439666 & 18439666 & 18439666 & 18439666 \\
    \texttt{product\_id} & 4 & 5 & 6 & 7 \\
    \texttt{product\_name} & Product4 & Product5 & Product6 & Product7 \\
    \texttt{product\_variant} & yellow & red & orange & grey \\
    \texttt{product\_brand} & Apple & Samsung & Coca-Cola & Audi \\
    \texttt{product\_description} & Description3 & Description4 & Description5 & Description6 \\
    \texttt{product\_image} & https://\allowbreak example.com/\allowbreak image3.jpg & https://\allowbreak example.com/\allowbreak image4.jpg & https://\allowbreak example.com/\allowbreak image5.jpg & https://\allowbreak example.com/\allowbreak image6.jpg \\
    \texttt{category\_id} & 4 & 5 & 6 & 7 \\
    \texttt{category\_name} & Categoria4 & Categoria5 & Categoria6 & Categoria7 \\
    \texttt{quantity\_available} & 82 & 57 & 64 & 1 \\
    \texttt{price} & 55.94 & 34.21 & 51.98 & 100000 \\
    \texttt{last\_restock\_date} & 2023-11-02 14:30:00Z & 2023-11-02 17:15:00Z & 2023-11-02 20:00:00Z & 2023-11-02 22:30:00Z \\
    \bottomrule
  \end{tabular}
  \par\vspace{4pt}
  {\footnotesize Fonte: Elaborado pelo autor (2026), a partir de \texttt{data/ecommerce/ecommerce\_join.csv}.}
\end{table}

\section{Linhas com \texttt{action == purchase}}

As oito linhas de compra são as mais completas do arquivo (34 das 39 colunas):
acrescentam aos grupos anteriores os dados de pedido, contraparte da venda, avaliação e
pagamento. As colunas de produto, categoria e estoque trazem os dados cadastrais do
produto negociado, tal como nas linhas de estoque (Tabela~\ref{tab:csv-stock-completa}),
e as colunas \texttt{trader\_*} trazem o endereço do vendedor. As
Tabelas~\ref{tab:csv-purchase-completa} e~\ref{tab:csv-purchase-completa-b} apresentam
os valores.

\begin{table}[H]
  \caption{Linhas \texttt{action == purchase} do \texttt{ecommerce\_join.csv}, todas as colunas preenchidas (linhas 1 a 4).}
  \label{tab:csv-purchase-completa}
  \centering
  \setlength{\tabcolsep}{3pt}%
  \renewcommand{\arraystretch}{1.15}%
  \begin{tabular}{l *{4}{>{\raggedright\arraybackslash}p{2.5cm}}}
    \toprule
    \textbf{Coluna} & \textbf{Linha 1} & \textbf{Linha 2} & \textbf{Linha 3} & \textbf{Linha 4} \\
    \midrule
    \texttt{action} & purchase & purchase & purchase & purchase \\
    \texttt{timestamp} & 2023-11-15 15:46:30Z & 2023-11-16 03:46:30Z & 2023-11-16 18:46:30Z & 2023-11-17 06:46:30Z \\
    \texttt{user\_id} & user1 & user2 & user3 & user4 \\
    \texttt{user\_name} & user1 & user2 & user3 & user4 \\
    \texttt{user\_email} & user1@\allowbreak example.com & user2@\allowbreak example.com & user3@\allowbreak example.com & user4@\allowbreak example.com \\
    \texttt{street} & Street 1, 42 & Street 2, 42 & Street 3, 42 & Street 4, 42 \\
    \texttt{neighborhood} & ABC & DEF & GHI & JKL \\
    \texttt{state} & CA & SC & RS & TO \\
    \texttt{country} & United States & Brasil & Brasil & Brasil \\
    \texttt{zip\_code} & 12345 & 84102149 & 18057313 & 19645442 \\
    \texttt{order\_number} & ORD1 & ORD2 & ORD3 & ORD4 \\
    \texttt{order\_date} & 2023-11-03T14:30:00-05:00 & 2023-11-03T12:30:00-03:00 & 2023-11-05T15:30:00-03:00 & 2023-11-07T16:20:00-03:00 \\
    \texttt{order\_status} & Completed & Completed & Completed & Completed \\
    \texttt{traded\_with} & user2 & user1 & user9 & user9 \\
    \texttt{trader\_street} & Street 2, 42 & Street 1, 42 & Street 9, 42 & Street 9, 42 \\
    \texttt{trader\_neighborhood} & DEF & ABC & YZ & YZ \\
    \texttt{trader\_state} & SC & CA & AM & AM \\
    \texttt{trader\_country} & Brasil & United States & Brasil & Brasil \\
    \texttt{trader\_zip\_code} & 84102149 & 12345 & 18439666 & 18439666 \\
    \texttt{comment} & Bad & Bom & Médio & Ótimo produto \\
    \texttt{rating} & 2 & 4 & 3 & 5 \\
    \texttt{payment\_method} & cash & cash & credit\_\allowbreak card & credit\_\allowbreak card \\
    \texttt{quantity} & 1 & 1 & 1 & 1 \\
    \texttt{product\_id} & 2 & 3 & 4 & 5 \\
    \texttt{product\_name} & Product2 & Product3 & Product4 & Product5 \\
    \texttt{product\_variant} & blue & green & yellow & red \\
    \texttt{product\_brand} & Microsoft & Google & Apple & Samsung \\
    \texttt{product\_description} & Description1 & Description2 & Description3 & Description4 \\
    \texttt{product\_image} & https://\allowbreak example.com/\allowbreak image1.jpg & https://\allowbreak example.com/\allowbreak image2.jpg & https://\allowbreak example.com/\allowbreak image3.jpg & https://\allowbreak example.com/\allowbreak image4.jpg \\
    \texttt{category\_id} & 2 & 3 & 4 & 5 \\
    \texttt{category\_name} & Categoria2 & Categoria3 & Categoria4 & Categoria5 \\
    \texttt{quantity\_available} & 2 & 1 & 82 & 57 \\
    \texttt{price} & 50.5 & 52.07 & 55.94 & 34.21 \\
    \texttt{last\_restock\_date} & 2023-11-02 03:30:00Z & 2023-11-02 06:45:00Z & 2023-11-02 14:30:00Z & 2023-11-02 17:15:00Z \\
    \bottomrule
  \end{tabular}
  \par\vspace{4pt}
  {\footnotesize Fonte: Elaborado pelo autor (2026), a partir de \texttt{data/ecommerce/ecommerce\_join.csv}.}
\end{table}

\begin{table}[H]
  \caption{Linhas \texttt{action == purchase} do \texttt{ecommerce\_join.csv}, todas as colunas preenchidas (linhas 5 a 8).}
  \label{tab:csv-purchase-completa-b}
  \centering
  \setlength{\tabcolsep}{3pt}%
  \renewcommand{\arraystretch}{1.15}%
  \begin{tabular}{l *{4}{>{\raggedright\arraybackslash}p{2.5cm}}}
    \toprule
    \textbf{Coluna} & \textbf{Linha 5} & \textbf{Linha 6} & \textbf{Linha 7} & \textbf{Linha 8} \\
    \midrule
    \texttt{action} & purchase & purchase & purchase & purchase \\
    \texttt{timestamp} & 2023-11-17 22:46:30Z & 2023-11-18 10:46:30Z & 2023-11-18 23:46:30Z & 2023-11-19 15:46:30Z \\
    \texttt{user\_id} & user5 & user6 & user7 & user8 \\
    \texttt{user\_name} & user5 & user6 & user7 & user8 \\
    \texttt{user\_email} & user5@\allowbreak example.com & user6@\allowbreak example.com & user7@\allowbreak example.com & user8@\allowbreak example.com \\
    \texttt{street} & Street 5, 42 & Street 6, 42 & Street 7, 42 & Street 8, 42 \\
    \texttt{neighborhood} & MNO & PQR & STU & VWX \\
    \texttt{state} & RJ & SP & SP & SE \\
    \texttt{country} & Brasil & Brasil & Brasil & Brasil \\
    \texttt{zip\_code} & 30097508 & 92083256 & 67065799 & 18439703 \\
    \texttt{order\_number} & ORD5 & ORD6 & ORD7 & ORD8 \\
    \texttt{order\_date} & 2023-11-09T12:30:00-03:00 & 2023-11-10T12:30:00-03:00 & 2023-11-11T12:30:00-03:00 & 2023-11-14T12:30:00-03:00 \\
    \texttt{order\_status} & Completed & Completed & Completed & Completed \\
    \texttt{traded\_with} & user9 & user9 & user8 & user2 \\
    \texttt{trader\_street} & Street 9, 42 & Street 9, 42 & Street 8, 42 & Street 2, 42 \\
    \texttt{trader\_neighborhood} & YZ & YZ & VWX & DEF \\
    \texttt{trader\_state} & AM & AM & SE & SC \\
    \texttt{trader\_country} & Brasil & Brasil & Brasil & Brasil \\
    \texttt{trader\_zip\_code} & 18439666 & 18439666 & 18439703 & 84102149 \\
    \texttt{comment} & Não gostei & --- & Reprovado & Bom, bonito e barato \\
    \texttt{rating} & 1 & 5 & 1 & 4 \\
    \texttt{payment\_method} & credit\_\allowbreak card & credit\_\allowbreak card & credit\_\allowbreak card & credit\_\allowbreak card \\
    \texttt{quantity} & 1 & 1 & 1 & 1 \\
    \texttt{product\_id} & 6 & 7 & 8 & 9 \\
    \texttt{product\_name} & Product6 & Product7 & Product8 & Product9 \\
    \texttt{product\_variant} & orange & grey & white & black \\
    \texttt{product\_brand} & Coca-Cola & Audi & Google & Pepsi \\
    \texttt{product\_description} & Description5 & Description6 & Description7 & Description8 \\
    \texttt{product\_image} & https://\allowbreak example.com/\allowbreak image5.jpg & https://\allowbreak example.com/\allowbreak image6.jpg & https://\allowbreak example.com/\allowbreak image7.jpg & https://\allowbreak example.com/\allowbreak image8.jpg \\
    \texttt{category\_id} & 6 & 7 & 8 & 9 \\
    \texttt{category\_name} & Categoria6 & Categoria7 & Categoria8 & Categoria9 \\
    \texttt{quantity\_available} & 64 & 1 & 60 & 67 \\
    \texttt{price} & 51.98 & 100000 & 46.79 & 32.78 \\
    \texttt{last\_restock\_date} & 2023-11-02 20:00:00Z & 2023-11-01T00:00:00-03:00 & 2023-11-02 12:00:00Z & 2023-11-02 09:15:00Z \\
    \bottomrule
  \end{tabular}
  \par\vspace{4pt}
  {\footnotesize Fonte: Elaborado pelo autor (2026), a partir de \texttt{data/ecommerce/ecommerce\_join.csv}.}
\end{table}

\section{Linhas com \texttt{action == select\_product}}

As oito linhas de seleção de produto preenchem apenas a identificação do evento, o
usuário e o grupo de seleção (7 das 39 colunas): o produto visualizado
(\texttt{selected\_product\_id}) e a quantidade de sugestões exibidas
(\texttt{suggested\_product\_count}). As Tabelas~\ref{tab:csv-select-completa}
e~\ref{tab:csv-select-completa-b} apresentam os valores.

\begin{table}[H]
  \caption{Linhas \texttt{action == select\_product} do \texttt{ecommerce\_join.csv}, todas as colunas preenchidas (linhas 1 a 4).}
  \label{tab:csv-select-completa}
  \centering
  \setlength{\tabcolsep}{3pt}%
  \renewcommand{\arraystretch}{1.15}%
  \begin{tabular}{l *{4}{>{\raggedright\arraybackslash}p{2.5cm}}}
    \toprule
    \textbf{Coluna} & \textbf{Linha 1} & \textbf{Linha 2} & \textbf{Linha 3} & \textbf{Linha 4} \\
    \midrule
    \texttt{action} & select\_\allowbreak product & select\_\allowbreak product & select\_\allowbreak product & select\_\allowbreak product \\
    \texttt{timestamp} & 2023-11-15 15:09:30Z & 2023-11-16 03:11:30Z & 2023-11-16 17:50:30Z & 2023-11-17 06:05:30Z \\
    \texttt{user\_id} & user1 & user2 & user3 & user4 \\
    \texttt{user\_name} & user1 & user2 & user3 & user4 \\
    \texttt{user\_email} & user1@\allowbreak example.com & user2@\allowbreak example.com & user3@\allowbreak example.com & user4@\allowbreak example.com \\
    \texttt{selected\_product\_id} & 11 & 12 & 22 & 14 \\
    \texttt{suggested\_product\_count} & 3 & 3 & 3 & 3 \\
    \bottomrule
  \end{tabular}
  \par\vspace{4pt}
  {\footnotesize Fonte: Elaborado pelo autor (2026), a partir de \texttt{data/ecommerce/ecommerce\_join.csv}.}
\end{table}

\begin{table}[H]
  \caption{Linhas \texttt{action == select\_product} do \texttt{ecommerce\_join.csv}, todas as colunas preenchidas (linhas 5 a 8).}
  \label{tab:csv-select-completa-b}
  \centering
  \setlength{\tabcolsep}{3pt}%
  \renewcommand{\arraystretch}{1.15}%
  \begin{tabular}{l *{4}{>{\raggedright\arraybackslash}p{2.5cm}}}
    \toprule
    \textbf{Coluna} & \textbf{Linha 5} & \textbf{Linha 6} & \textbf{Linha 7} & \textbf{Linha 8} \\
    \midrule
    \texttt{action} & select\_\allowbreak product & select\_\allowbreak product & select\_\allowbreak product & select\_\allowbreak product \\
    \texttt{timestamp} & 2023-11-17 21:56:30Z & 2023-11-18 10:09:30Z & 2023-11-18 23:20:30Z & 2023-11-19 15:18:30Z \\
    \texttt{user\_id} & user5 & user6 & user7 & user8 \\
    \texttt{user\_name} & user5 & user6 & user7 & user8 \\
    \texttt{user\_email} & user5@\allowbreak example.com & user6@\allowbreak example.com & user7@\allowbreak example.com & user8@\allowbreak example.com \\
    \texttt{selected\_product\_id} & 14 & 21 & 33 & 17 \\
    \texttt{suggested\_product\_count} & 3 & 3 & 3 & 3 \\
    \bottomrule
  \end{tabular}
  \par\vspace{4pt}
  {\footnotesize Fonte: Elaborado pelo autor (2026), a partir de \texttt{data/ecommerce/ecommerce\_join.csv}.}
\end{table}

\section{Linhas com \texttt{action == add\_to\_cart}}

As oito linhas de adição ao carrinho preenchem a identificação do evento, o usuário e
o grupo de carrinho (8 das 39 colunas). O identificador do carrinho
(\texttt{shopping\_cart\_id}) segue o padrão
\texttt{user:$\langle$user\_id$\rangle$:cart}. As Tabelas~\ref{tab:csv-cart-completa}
e~\ref{tab:csv-cart-completa-b} apresentam os valores.

\begin{table}[H]
  \caption{Linhas \texttt{action == add\_to\_cart} do \texttt{ecommerce\_join.csv}, todas as colunas preenchidas (linhas 1 a 4).}
  \label{tab:csv-cart-completa}
  \centering
  \setlength{\tabcolsep}{3pt}%
  \renewcommand{\arraystretch}{1.15}%
  \begin{tabular}{l *{4}{>{\raggedright\arraybackslash}p{2.5cm}}}
    \toprule
    \textbf{Coluna} & \textbf{Linha 1} & \textbf{Linha 2} & \textbf{Linha 3} & \textbf{Linha 4} \\
    \midrule
    \texttt{action} & add\_\allowbreak to\_\allowbreak cart & add\_\allowbreak to\_\allowbreak cart & add\_\allowbreak to\_\allowbreak cart & add\_\allowbreak to\_\allowbreak cart \\
    \texttt{timestamp} & 2023-11-15 15:17:30Z & 2023-11-16 03:27:30Z & 2023-11-16 18:08:30Z & 2023-11-17 06:14:30Z \\
    \texttt{user\_id} & user1 & user2 & user3 & user4 \\
    \texttt{user\_name} & user1 & user2 & user3 & user4 \\
    \texttt{user\_email} & user1@\allowbreak example.com & user2@\allowbreak example.com & user3@\allowbreak example.com & user4@\allowbreak example.com \\
    \texttt{shopping\_cart\_id} & {\footnotesize user:\allowbreak 7fd7eee7-\allowbreak 4675-\allowbreak 4f56-\allowbreak ad81-\allowbreak 108da6cecebc:\allowbreak cart} & {\footnotesize user:\allowbreak f246fb91-\allowbreak e47d-\allowbreak 44e8-\allowbreak a015-\allowbreak ce5dfd2cb8d1:\allowbreak cart} & {\footnotesize user:\allowbreak 09b3458a-\allowbreak 9d1d-\allowbreak 4708-\allowbreak 8d2d-\allowbreak 3eb5fadbac9b:\allowbreak cart} & {\footnotesize user:\allowbreak 1fd3d09a-\allowbreak a310-\allowbreak 4a46-\allowbreak bfdf-\allowbreak 37f2f9935a8c:\allowbreak cart} \\
    \texttt{cart\_product\_id} & 3 & 2 & 2 & 2 \\
    \texttt{cart\_quantity} & 1 & 1 & 1 & 1 \\
    \bottomrule
  \end{tabular}
  \par\vspace{4pt}
  {\footnotesize Fonte: Elaborado pelo autor (2026), a partir de \texttt{data/ecommerce/ecommerce\_join.csv}.}
\end{table}

\begin{table}[H]
  \caption{Linhas \texttt{action == add\_to\_cart} do \texttt{ecommerce\_join.csv}, todas as colunas preenchidas (linhas 5 a 8).}
  \label{tab:csv-cart-completa-b}
  \centering
  \setlength{\tabcolsep}{3pt}%
  \renewcommand{\arraystretch}{1.15}%
  \begin{tabular}{l *{4}{>{\raggedright\arraybackslash}p{2.5cm}}}
    \toprule
    \textbf{Coluna} & \textbf{Linha 5} & \textbf{Linha 6} & \textbf{Linha 7} & \textbf{Linha 8} \\
    \midrule
    \texttt{action} & add\_\allowbreak to\_\allowbreak cart & add\_\allowbreak to\_\allowbreak cart & add\_\allowbreak to\_\allowbreak cart & add\_\allowbreak to\_\allowbreak cart \\
    \texttt{timestamp} & 2023-11-17 22:11:30Z & 2023-11-18 10:21:30Z & 2023-11-18 23:34:30Z & 2023-11-19 15:29:30Z \\
    \texttt{user\_id} & user5 & user6 & user7 & user8 \\
    \texttt{user\_name} & user5 & user6 & user7 & user8 \\
    \texttt{user\_email} & user5@\allowbreak example.com & user6@\allowbreak example.com & user7@\allowbreak example.com & user8@\allowbreak example.com \\
    \texttt{shopping\_cart\_id} & {\footnotesize user:\allowbreak 0331463b-\allowbreak 7e62-\allowbreak 44ed-\allowbreak 945f-\allowbreak 6b3c8195245b:\allowbreak cart} & {\footnotesize user:\allowbreak f4857da9-\allowbreak c991-\allowbreak 4bc6-\allowbreak b0a9-\allowbreak a4277c005956:\allowbreak cart} & {\footnotesize user:\allowbreak a87df805-\allowbreak ad1b-\allowbreak 4ae9-\allowbreak b5ad-\allowbreak cbb64ce9e475:\allowbreak cart} & {\footnotesize user:\allowbreak c11a1ef2-\allowbreak c7c0-\allowbreak 4a47-\allowbreak 8128-\allowbreak 16dc3da9d72b:\allowbreak cart} \\
    \texttt{cart\_product\_id} & 2 & 2 & 2 & 2 \\
    \texttt{cart\_quantity} & 1 & 1 & 1 & 1 \\
    \bottomrule
  \end{tabular}
  \par\vspace{4pt}
  {\footnotesize Fonte: Elaborado pelo autor (2026), a partir de \texttt{data/ecommerce/ecommerce\_join.csv}.}
\end{table}
